\clearpage
\section{Processor-State Save and Restore}\label{page:processor-state-save-and-restore}

SAVE and RESTORE use a CPUID-defined memory image whose base is 4 KiB aligned.

\subsection{Save-Area Layout}
Software obtains every extension component's offset and maximum size from CPUID
\texttt{SAVE\_AREA\_LAYOUT}. Each component starts on a 64-byte boundary.

\begin{manuallistedformatdiagram}{SAVE/RESTORE Base Header}
\manualformatrowrange{header[63:0]}{63}{0}{%
\manualformatreserved{reserved}{12}
\manualformatfield{STATUS}{16}
\manualformatfield{FLAGS}{16}
\manualformatreserved{reserved}{10}
\manualformatfield{5}{1}
\manualformatfield{4}{1}
\manualformatfield{3}{1}
\manualformatfield{2}{1}
\manualformatfield{1}{1}
\manualformatfield{0}{1}
\manualformatfield{FMT}{4}
}
\manualformatrowrange{bitmap[63:0]}{63}{0}{%
\manualformatfield{state block bitmap}{64}
}
\end{manuallistedformatdiagram}
{\small The labels \texttt{5} through \texttt{0} are the six \texttt{GS\_VALID} bits. SAVE writes
\texttt{GS\_VALID=0x3f} and \texttt{FMT=0}.\par}
\begin{manuallistedstructlayout}
  {SAVE/RESTORE Fixed Base Block}
  {8}
  {offset from the 4 KiB-aligned save-area base}
  {offsets increase downward}
\manualstructrow{0x000--0x00f}{%
  \manualstructheaderfield{4}{BASE HEADER\\\texttt{0x000} \textperiodcentered\ 8 bytes}
  \manualstructheaderfield{4}{STATE BLOCK BITMAP\\\texttt{0x008} \textperiodcentered\ 8 bytes}
}
\manualstructrow{0x010--0x04f}{%
  \manualstructslotseries{R}{0}{8}
}
\manualstructrow{0x050--0x08f}{%
  \manualstructslotseries{R}{8}{8}
}
\manualstructrow{0x090--0x0bf}{%
  \manualstructslotseries{GS}{0}{6}
  \manualstructannotationfield{2}{RESTORE uses slots\\selected by \texttt{GS\_VALID}}
}
\manualstructtallrow{0x0c0+}{%
  \manualstructextensionfield{8}{EXTENSION STATE COMPONENTS\\CPUID-defined slots \textperiodcentered\ each component starts on a 64-byte boundary}
}
\manualstructfooter
  {Each \texttt{Rn} and \texttt{GSn} cell is one 8-byte slot.}
  {\texttt{Rn}=\texttt{0x010}+8n; \texttt{GSn}=\texttt{0x090}+8n.}
\end{manuallistedstructlayout}


\subsection{Fixed Architectural State}
The fixed block contains the base header, state-block bitmap, R0 through R15, and GS0 through GS5 images.

The header carries FLAGS, STATUS, format, and GS-validity state. \texttt{GS\_VALID} selects which GS images are
meaningful.

\subsection{Extension Components and Clean State}
An extension component is architecturally clean when its state is known to equal its component-defined initial
state. Reset and RESTORE of a clear component bit make that component clean. A committed write that can make a
component differ from its initial state makes it modified; RESTORE of a non-initial component image also makes it
modified. An implementation may mark a component clean whenever it establishes that the complete component again
equals its initial state.

When floating-point state is enabled, its component contains F0 through F15, FFLAGS, and FSTATUS in the layout
reported in the CPUID chapter. The floating-point component's initial state is positive zero in F0 through F15 and
zero in FFLAGS and FSTATUS.

\subsection{SAVE Responsibilities}
SAVE writes \texttt{FMT=0}, sets \texttt{GS\_VALID=0x3f}, and records GS0 through GS5. SAVE sets bitmap bits only
for components described by its CPUID \texttt{SAVE\_AREA\_LAYOUT} leaf.

SAVE may omit an architecturally clean extension component by clearing its bitmap bit. SAVE may clear a bitmap bit
only for a clean component. SAVE does not change clean or modified state.

\subsection{RESTORE Responsibilities}
RESTORE replaces each GSn whose corresponding \texttt{GS\_VALID} bit is set and preserves each GSn whose bit is
clear. It validates every selected segment image before changing architectural state.

RESTORE uses the state-block bitmap as authoritative when reconstructing extension state. It accepts only the
defined format and only bitmap bits that name CPUID-described components for that format. RESTORE raises
\texttt{INVALID\_CONTROL\_STATE} before changing architectural state when the format is unrecognized or a set
bitmap bit does not name a CPUID-described component for that format. RESTORE of a clear, described component bit
installs that component's initial state.

User-mode RESTORE ignores saved supervisor-only STATUS state. Supervisor RESTORE validates the saved STATUS image by
the ordinary privileged write rules before making it visible.

The normative instruction entries are
\hyperref[instr:save]{\texttt{SAVE}} and
\hyperref[instr:restore]{\texttt{RESTORE}}.\par
